% ============================================================
% СХЕМА ЭЛЕКТРИЧЕСКАЯ СОЕДИНЕНИЙ РЯДОВ ЗАЖИМОВ
% ============================================================

{\color{unidarkgreen}\section[(обязательное) Схема электрическая соединений рядов зажимов]{(обязательное)\\Схема электрическая соединений рядов зажимов}\label{app:a41}}
\color{black}

% Часть 1 – размер 0.9 (если нужно, можете изменить на 0.95)
\begin{figure}[H]
  \centering
  \includegraphics[width=0.9\textwidth, height=0.9\textheight, keepaspectratio, trim=5pt 5pt 5pt 5pt, clip]{\apppath/Приложение. Схема Э4.1/_latex/img/Fig_1.pdf}
  \caption{Схема электрическая соединений рядов зажимов (часть 1)}
  \label{appvo:fig_1}
\end{figure}

% Часть 2
\begin{figure}[H]
  \centering
  \includegraphics[width=0.95\textwidth, height=0.95\textheight, keepaspectratio, trim=5pt 5pt 5pt 5pt, clip]{\apppath/Приложение. Схема Э4.1/_latex/img/Fig_2.pdf}
  \caption{Схема электрическая соединений рядов зажимов (часть 2)}
  \label{appvo:fig_2}
\end{figure}

% Часть 3
\begin{figure}[H]
  \centering
  \includegraphics[width=0.95\textwidth, height=0.95\textheight, keepaspectratio, trim=5pt 5pt 5pt 5pt, clip]{\apppath/Приложение. Схема Э4.1/_latex/img/Fig_3.pdf}
  \caption{Схема электрическая соединений рядов зажимов (часть 3)}
  \label{appvo:fig_3}
\end{figure}

% Часть 4
\begin{figure}[H]
  \centering
  \includegraphics[width=0.95\textwidth, height=0.95\textheight, keepaspectratio, trim=5pt 5pt 5pt 5pt, clip]{\apppath/Приложение. Схема Э4.1/_latex/img/Fig_4.pdf}
  \caption{Схема электрическая соединений рядов зажимов (часть 4)}
  \label{appvo:fig_4}
\end{figure}

\clearpage